\documentclass[11pt,a4paper]{article}
\usepackage[margin=24mm]{geometry}
\usepackage{amsmath,amssymb,amsthm}
\setlength{\emergencystretch}{2em}
\newtheorem{theorem}{Theorem}
\newtheorem{corollary}[theorem]{Corollary}
\newcommand{\gen}[1]{\langle#1\rangle}
\title{Completion of the singleton and unique-crossing branches}
\author{Human Compression project: L044 continuation}
\date{}
\begin{document}\maketitle
Let $G\le\operatorname{Alt}(\Omega)$ be transitive with a connected
string C-marking $\rho_0,\ldots,\rho_{r-1}$, and suppose every
maximal ambient omission is intransitive. Write $n=|\Omega|\ge12$.
Assume the alternating rank bound at every strictly smaller degree
at least twelve. All complete acquired primitive, anchored-imprimitive,
small-degree and attachment inputs are retained with their actual
scope; this induction premise does not presume the theorem at degree $n$.
Eligibility means crossing two actual ambient omission orbits.

\begin{theorem}[A singleton at an end]\label{singleton}
If $G_0$ has the two orbits $\{a\}$ and
$\Omega\setminus\{a\}$, then
$r\le\lfloor(n-1)/2\rfloor$.
The same conclusion holds with the string reversed.
\end{theorem}
\begin{proof}
Suppose the bound fails. The complete critical odd-tail reduction
gives $n=2r\ge14$ and the actual natural
$G_0=\operatorname{Alt}(\Omega\setminus\{a\})$.
The complete critical-head theorem excludes either actual supported
six-point $A_5$ end from source Table~2(2)--(3), applying full
string reversal at the other end. Its small boundary and the
unextendable orientation are retained in that theorem's proof.

The complete singleton global-pairs theorem supplies a further
colour $j\ge2$ with exactly one ambient-eligible edge. Choose
the least such index. Every colour $2,\ldots,j-1$ is paired:
it has at least one eligible edge, since otherwise adjoining its
mark could not make its intransitive omission transitive, and
minimality excludes exactly one.

By the actual singleton-prefix lemma the two $G_j$-orbits are
\[
 C=a^{G_{<j}},\qquad A=\Omega\setminus C,
\]
the suffix $G_{>j}$ fixes $C$ pointwise, and every shared
label at this cut is below $j$. Strict growth gives $|C|\ge j+1$.

First suppose $A=\{d\}$ is a singleton. If $j<r-1$, the
complete interior singleton-orbit theorem gives the required bound,
contrary to our assumption. If $j=r-1$, reverse the whole string
and apply the critical odd-tail reduction to this actual endpoint
omission. It gives
$G_{r-1}=\operatorname{Alt}(\Omega\setminus\{d\})$.
The minimal choice of $j$ supplies every paired colour
$2,\ldots,r-2$. These are precisely the hypotheses of the complete
two-singleton theorem, which excludes the configuration. Thus
$A$ is nontrivial from now on.

If no label is shared across the cut, each nonidentity mark other
than $\rho_j$ acts on exactly one of $C,A$. Within either of
the two string intervals the side cannot change: otherwise two
adjacent marks would have disjoint supports and commute, disconnecting
the full string diagram. The prefix acts on $C$, since $\rho_0$
moves $a$. Nontriviality of $A$ forces $j<r-1$ and makes the
suffix act on $A$. Thus these are actual supported transitive
prefix and suffix groups, their transitivity following from the
two actual omission orbits. The complete supported-interval
alternative gives the rank bound or an actual exceptional end,
both already excluded.

Consequently a first shared label $g<j$ exists. The complete
singleton-prefix support theorem gives
\[
             g\ge3,\qquad
             Y=\operatorname{supp}(G_{<g})\subseteq C,
             \qquad |Y|\le2g+1.                         \tag{1}
\]
Here the support inclusion is actual: every earlier unshared mark
acts on the same side, namely $C$, because the prefix starts at
$\rho_0$. The small cases $g=1,2$ have already been excluded
in that theorem, including its full critical-dihedral argument
at the original reversed cut.

If $G_{<g}$ is transitive on $Y$, use the full transitive-tail
alternative after reversing the entire string. Its crossing is
$i=r-1-j$ and its last shared label is $h=r-1-g$, so
$i<h<r-1$. Its supported tail past $h$ is exactly the actual
$G_{<g}$ on $Y\subseteq C$. If $i>0$, the preceding subgroup
in this orientation is the original $G_{>j}$, which fixes $C$
pointwise. Thus the chosen-orientation condition is proved,
not inferred from a projected marking. If $i=0$, the endpoint
form applies. In both cases the theorem gives the rank bound
or an actual exceptional end, again impossible.

We are left with $G_{<g}$ intransitive on its support.
At $g=3$, (1) contradicts the complete third-prefix hand proof,
which requires at least eight moved points. For every $g\ge4$,
put $H=G_{\{2,\ldots,g-1\}}$. All these colours are paired
by $g<j$ and the least choice of $j$. The complete small-support
budget therefore applies with its exact actual hypotheses.
If $f$ counts the $H$-fixed points in $Y$ and $c$ its nontrivial
orbits, its exhaustive table is
\[
 (f,c,|Y|)\in
 \{(2,2,2g),(2,2,2g+1),(3,2,2g+1),(2,3,2g+1)\}.
\]
The complete two-orbit theorem excludes the first two entries;
the third-fixed-point theorem excludes the third; the selected
ambient confinement and three-orbit theorem exclude the fourth.
All hypotheses concern the actual prefix and actual ambient
eligible edges. No pure alternating prefix kernel or projected
intersection property is used.

This exhausts every later cut and every possible first shared
prefix. The assumed rank failure is impossible.
\end{proof}

\begin{corollary}[The unique-crossing branch]\label{unique}
Under the opening hypotheses, if at least one colour has exactly
one ambient-eligible edge, then
$r\le\lfloor(n-1)/2\rfloor$.
\end{corollary}
\begin{proof}
Suppose the bound fails and choose such a crossing of index $i$.
Its omission has exactly two orbits. If either is a singleton,
the interior singleton-orbit theorem or Theorem~\ref{singleton},
with full reversal when necessary, excludes it.

Suppose both orbits are nontrivial. If no label is shared, the
same connected-interval argument used above puts all marks on
opposite sides of the crossing. The crossing is interior, since
at an endpoint its only remaining interval could act on just
one orbit, leaving the other a singleton. The supported-interval
alternative therefore applies and gives an exceptional actual
end, excluded by the complete critical-head theorem.

If some label is shared, reverse the full string if necessary
so that at least one unique crossing has a shared label to its
right. Choose the maximal index among all such crossings and
apply the complete maximal-crossing alternative. It gives the
rank bound, an actual exceptional end, or the terminal singleton
case. Its retained proof obtains this last case at an actual
later \emph{unique} crossing $j=r-1$; hence that endpoint omission
has exactly the singleton and its complement as its two orbits.
This is the full two-orbit hypothesis needed by
Theorem~\ref{singleton}, not merely an arbitrary fixed point of
an omission. Each alternative is now excluded.
\end{proof}

\paragraph{Attribution and principal-source boundary.}
The cut strategy and rank target are those of
Cameron--Fernandes--Leemans--Mixer, \emph{Highest rank of a polytope
for $A_n$}, arXiv:1605.09173v1, Section~5. This assembly retains
all complete acquired proofs, including explicit smaller-degree
induction and their required finite inputs. The new uniform prefix
orbit exclusions supply the missing singleton continuation. The
separately reviewed fourth-prefix GAP certificate is unnecessary
as an additional premise here, since the uniform budget covers
every $g\ge4$.

The conclusion covers the all-intransitive-omission branch with
at least one unique crossing. The all-paired equality case remains
separate, as do general transitive imprimitive maximal omissions,
the exact base maxima, lower-existence bindings and the complete
principal-theorem assembly. No whole L044 recovery, conjecture,
measurement or new formal declaration is claimed.
\end{document}
